% ============================================================================
% CAPÍTULO 2 - DADOS E PREPARAÇÃO
% ============================================================================

\chapter{Dados e Preparação}\label{cap:dados}

Este capítulo descreve as quatro fontes de dados do projeto: os protocolos clínicos reais (Linhas de Cuidado), os dados sintéticos de pacientes, o banco de dados PostgreSQL e o dataset MedQuAD para fine-tuning.

\section{Linhas de Cuidado (Protocolos Clínicos)}

A base de conhecimento clínico do sistema é composta por \textbf{14 Linhas de Cuidado} --- protocolos clínicos profissionais em português brasileiro, cobrindo condições que vão de hipertensão e diabetes até demência e polifarmácia. Estes protocolos são documentos reais de prática clínica, com estrutura padronizada que inclui seções como Considerações Iniciais, Diagnóstico, Monitoramento, Encaminhamento e Indicadores.

\subsection{Protocolos Disponíveis}

A Tabela~\ref{tab:linhas_cuidado} lista os 14 protocolos utilizados como base de conhecimento para o sistema RAG.

\begin{table}[H]
\centering
\caption{14 Linhas de Cuidado indexadas no sistema.}
\label{tab:linhas_cuidado}
\rowcolors{2}{fiap-light}{white}
\begin{tabularx}{\textwidth}{c l l X}
\toprule
\textbf{\#} & \textbf{Pathway} & \textbf{Domínio} & \textbf{Exemplos de Conteúdo} \\
\midrule
1  & Hipertensão          & Cardiovascular    & Classificação PA, MAPA, anti-hipertensivos \\
2  & Diabetes             & Endócrino         & DM1/DM2, HbA1C, insulinização \\
3  & Asma                 & Respiratório      & Classificação gravidade, step therapy \\
4  & Depressão            & Psiquiátrico      & PHQ-9, ISRS, psicoterapia \\
5  & Ansiedade            & Psiquiátrico      & GAD-7, benzodiazepínicos, TCC \\
6  & DPOC                 & Respiratório      & GOLD, espirometria, oxigenoterapia \\
7  & Demência             & Neurológico       & Mini-Mental, anticolinesterásicos \\
8  & Dispepsia            & Gastroenterologia & H. pylori, IBP, endoscopia \\
9  & Hipotireoidismo      & Endócrino         & TSH, levotiroxina, seguimento \\
10 & Insuf.\ Cardíaca     & Cardiovascular    & NYHA, IECA, betabloqueadores \\
11 & Obesidade            & Metabólico        & IMC, dietoterapia, cirurgia bariátrica \\
12 & Pac.\ Cirúrgico      & Cirúrgico         & Pré-operatório, risco cirúrgico \\
13 & Polifarmácia         & Geriatria         & Reconciliação medicamentosa, Beers \\
14 & Tabagismo            & Preventivo        & Fagerström, TRN, bupropiona \\
\bottomrule
\end{tabularx}
\end{table}

\subsection{Extração e Limpeza dos Protocolos}

Os 14 protocolos foram disponibilizados originalmente em formato PDF. Para viabilizar a indexação no sistema RAG, foi realizado um processo de extração e limpeza:

\begin{enumerate}
    \item \textbf{Extração de texto:} Uso da biblioteca PyMuPDF (fitz) para extrair o conteúdo textual de cada página dos PDFs;
    \item \textbf{Identificação de seções:} Detecção automática de seções principais por meio de heurística baseada em texto em caixa alta (\textit{ALL CAPS}), seguindo o padrão estrutural dos documentos;
    \item \textbf{Limpeza:} Remoção de referências ao sistema de origem e caracteres indesejados;
    \item \textbf{Saída em Markdown:} Geração de arquivos \texttt{.md} com hierarquia de seções preservada (\texttt{\#\#} para seções, \texttt{\#\#\#} para subseções), facilitando o chunking posterior.
\end{enumerate}

O resultado são 14 arquivos Markdown limpos, totalizando aproximadamente 448~KB de conteúdo clínico estruturado, armazenados em \texttt{data/linhas\_de\_cuidado/extracted/}. Estes arquivos são versionados no repositório e servem como entrada direta para o pipeline de RAG (Capítulo~4).

\begin{infobox}[Nota sobre a Extração]
O script de extração de PDFs é um utilitário auxiliar executado uma única vez fora do pipeline principal. Os arquivos Markdown resultantes já estão incluídos no repositório, portanto \textbf{não é necessário reexecutar} a extração --- basta utilizar os arquivos \texttt{.md} já disponíveis.
\end{infobox}

\subsection{Protocolos Complementares de Emergência}

Para cobrir cenários de \textbf{emergência} não presentes nas 14 Linhas de Cuidado (que focam em condições crônicas), o sistema inclui 5 protocolos sintéticos adicionais. Estes protocolos foram elaborados com base em diretrizes reconhecidas (AHA, Surviving Sepsis Campaign, SBPT, IDSA) e seguem a mesma estrutura dos protocolos extraídos.

\begin{table}[H]
\centering
\caption{Protocolos de emergência sintéticos complementares.}
\label{tab:protocolos_emergencia}
\rowcolors{2}{fiap-light}{white}
\begin{tabularx}{\textwidth}{c l l X}
\toprule
\textbf{\#} & \textbf{Protocolo} & \textbf{CID} & \textbf{Conteúdo Principal} \\
\midrule
1 & IAM (Infarto Agudo do Miocárdio) & I21 & MONABCH, reperfusão, tempo porta-balão \\
2 & AVC (Acidente Vascular Cerebral) & I64 & NIHSS, trombólise rtPA, janela terapêutica \\
3 & Sepse e Choque Séptico & A41 & Hour-1 Bundle, qSOFA, ressuscitação \\
4 & Pneumonia (PAC) & J18 & CURB-65, antibioticoterapia empírica \\
5 & ITU (Infecção do Trato Urinário) & N39 & Cistite vs pielonefrite, bacteriúria \\
\bottomrule
\end{tabularx}
\end{table}

Cada protocolo sintético inclui as seções: Considerações Iniciais, Diagnóstico, Manejo Inicial, Monitoramento e Indicadores de Qualidade. Os arquivos são gerados automaticamente pelo script \texttt{01\_gerar\_dados\_sinteticos} e salvos em formato Markdown no diretório \texttt{data/synthetic/protocolos\_complementares/}, prontos para indexação no ChromaDB junto com as Linhas de Cuidado.

\begin{warningbox}[Protocolos Sintéticos]
Os protocolos de emergência são \textbf{sintéticos}, criados para fins educacionais e de demonstração. Em ambiente de produção, devem ser substituídos por protocolos institucionais validados pela equipe médica.
\end{warningbox}

\section{Dados Sintéticos de Pacientes}

Para simular um ambiente hospitalar realista, o sistema utiliza dados de pacientes inteiramente fictícios, gerados com coerência clínica. A geração é realizada pelo primeiro script do pipeline (\texttt{01\_gerar\_dados\_sinteticos}), que invoca o módulo \texttt{synthetic\_generator} localizado em \texttt{src/data/}.

\subsection{Geração de Pacientes}

Foram gerados \textbf{80 pacientes fictícios} utilizando a biblioteca Faker (localização \texttt{pt\_BR}), com sementes fixas (\texttt{seed=42}) para garantir a reprodutibilidade. Os dados pessoais incluem:

\begin{table}[H]
\centering
\caption{Campos do cadastro de pacientes sintéticos.}
\label{tab:campos_pacientes}
\rowcolors{2}{fiap-light}{white}
\begin{tabularx}{\textwidth}{l l X}
\toprule
\textbf{Campo} & \textbf{Exemplo} & \textbf{Geração} \\
\midrule
nome              & ``Maria Santos de Oliveira'' & \texttt{Faker.name()} \\
cpf               & ``123.456.789-00''           & \texttt{Faker.cpf()} \\
data\_nascimento  & 1958-03-15                   & \texttt{Faker.date\_of\_birth()} \\
sexo              & ``F''                        & Distribuição 55\%F / 45\%M \\
telefone          & ``(11) 98765-4321''          & \texttt{Faker.phone\_number()} \\
endereço          & ``Rua das Flores, 123''      & \texttt{Faker.address()} \\
comorbidades      & [``hipertensão'', ``DM2'']   & Distribuição probabilística \\
alergias          & [``dipirona'']               & Lista aleatória \\
medicamentos\_uso & [``losartana 50mg'']         & Coerente com comorbidades \\
\bottomrule
\end{tabularx}
\end{table}

\subsection{Coerência Clínica}

Um diferencial importante da geração de dados é a \textbf{coerência clínica} entre comorbidades, medicamentos, exames e consultas. O gerador implementa uma base de conhecimento médico com 9 condições alinhadas às Linhas de Cuidado:

\begin{table}[H]
\centering
\caption{Distribuição de condições nos dados sintéticos.}
\label{tab:distribuicao_condicoes}
\rowcolors{2}{fiap-light}{white}
\begin{tabularx}{\textwidth}{l c X}
\toprule
\textbf{Condição} & \textbf{Prevalência} & \textbf{Exemplos de Medicamentos Associados} \\
\midrule
Hipertensão         & 45\% & Losartana 50mg, enalapril 10mg, anlodipino 5mg \\
Diabetes tipo 2     & 30\% & Metformina 850mg, glibenclamida 5mg, insulina NPH \\
Obesidade           & 25\% & Sibutramina 10mg, orlistate 120mg \\
Depressão           & 20\% & Sertralina 50mg, fluoxetina 20mg, escitalopram 10mg \\
Ansiedade           & 15\% & Escitalopram 10mg, alprazolam 0,5mg \\
Asma                & 12\% & Salbutamol spray, budesonida 200mcg \\
Hipotireoidismo     & 10\% & Levotiroxina 50mcg, levotiroxina 75mcg \\
DPOC                & 8\%  & Tiotrópio 18mcg, formoterol + budesonida \\
Insuf.\ Cardíaca    & 5\%  & Enalapril 10mg, furosemida 40mg, carvedilol 12,5mg \\
\bottomrule
\end{tabularx}
\end{table}

Muitos pacientes possuem \textbf{múltiplas comorbidades} (por exemplo, hipertensão + diabetes + obesidade), o que é clinicamente realista e cria cenários ricos para o assistente consultar múltiplos protocolos simultaneamente.

\subsection{Dados Clínicos Gerados}

Para cada paciente, o gerador cria dados clínicos coerentes com suas comorbidades:

\begin{table}[H]
\centering
\caption{Volume de dados clínicos gerados.}
\label{tab:volume_dados}
\rowcolors{2}{fiap-light}{white}
\begin{tabular}{l r l}
\toprule
\textbf{Tipo} & \textbf{Total Gerado} & \textbf{Arquivo de Saída} \\
\midrule
Pacientes     & 80   & \texttt{pacientes.json} \\
Exames        & 199  & \texttt{exames.json} \\
Prontuários   & 297  & \texttt{consultas.json} \\
Receitas      & 163  & \texttt{receitas.json} \\
\bottomrule
\end{tabular}
\end{table}

Exemplos de coerência clínica implementada:
\begin{itemize}
    \item Paciente \textbf{diabético} $\rightarrow$ exames de HbA1C (valores entre 5,5--12\%), glicemia de jejum (80--250 mg/dL), creatinina;
    \item Paciente \textbf{hipertenso} $\rightarrow$ MAPA 24h com valores sistólico/diastólico realistas, ecocardiograma;
    \item Paciente \textbf{asmático} $\rightarrow$ espirometria com VEF1/CVF e classificação de gravidade;
    \item Paciente com \textbf{depressão} $\rightarrow$ avaliação PHQ-9 com escore de gravidade e notas clínicas compatíveis.
\end{itemize}

\section{Banco de Dados PostgreSQL}

Os dados sintéticos são armazenados em um banco PostgreSQL 16, executado em container Docker. O setup é automatizado pelo \texttt{scripts/02\_setup\_postgres.py}.

\subsection{Schema do Banco}

O banco possui 4 tabelas, modeladas com SQLAlchemy ORM:

\begin{lstlisting}[style=sql, caption={Schema do banco de dados.}]
-- Pacientes
CREATE TABLE pacientes (
    id UUID PRIMARY KEY,
    nome VARCHAR(200) NOT NULL,
    cpf VARCHAR(14) UNIQUE NOT NULL,
    data_nascimento DATE NOT NULL,
    sexo VARCHAR(1) NOT NULL,
    telefone VARCHAR(20),
    endereco TEXT,
    comorbidades TEXT[],
    alergias TEXT[],
    medicamentos_uso TEXT[]
);

-- Exames clinicos
CREATE TABLE exames (
    id UUID PRIMARY KEY,
    paciente_id UUID REFERENCES pacientes(id),
    tipo VARCHAR(100) NOT NULL,
    data DATE NOT NULL,
    resultados JSONB,       -- Resultados estruturados
    interpretacao TEXT,
    status VARCHAR(20)      -- pendente, realizado, cancelado
);

-- Prontuarios (consultas medicas)
CREATE TABLE prontuarios (
    id UUID PRIMARY KEY,
    paciente_id UUID REFERENCES pacientes(id),
    data DATE NOT NULL,
    diagnostico_cid VARCHAR(10),  -- Codigo CID-10
    diagnostico_texto TEXT,
    observacoes TEXT
);

-- Receitas medicas
CREATE TABLE receitas (
    id UUID PRIMARY KEY,
    paciente_id UUID REFERENCES pacientes(id),
    data DATE NOT NULL,
    condicao TEXT,
    medicamentos TEXT[],
    validade_dias INTEGER,
    observacoes TEXT
);
\end{lstlisting}

\subsection{Características Técnicas}

\begin{itemize}
    \item \textbf{UUIDs como chaves primárias:} Evita colisões e permite geração distribuída;
    \item \textbf{JSONB para resultados de exames:} Permite armazenar resultados com estrutura variável (diferentes tipos de exames geram diferentes campos), mantendo a capacidade de consulta;
    \item \textbf{ARRAY para listas:} Comorbidades, alergias e medicamentos são armazenados como arrays nativos do PostgreSQL, facilitando consultas com operador \texttt{ANY};
    \item \textbf{CID-10:} Diagnósticos são codificados no padrão \acrfull{cid}, permitindo consultas padronizadas;
    \item \textbf{Consultas parametrizadas:} O módulo \texttt{src/database/queries.py} utiliza \texttt{SQLAlchemy text()} com parâmetros nomeados para prevenir injeção SQL.
\end{itemize}

\subsection{Infraestrutura Docker}

O script \texttt{02\_setup\_postgres.py} automatiza todo o processo:

\begin{enumerate}
    \item Verifica se o Docker Desktop está em execução;
    \item Cria ou inicia o container \texttt{medical\_assistant\_db} (PostgreSQL 16 Alpine);
    \item Cria as 4 tabelas via \texttt{SQLAlchemy Base.metadata.create\_all()};
    \item Carrega os 4 arquivos JSON gerados pelo script~01;
    \item Popula o banco de dados com os dados sintéticos;
    \item Executa validação final (contagem de registros por tabela).
\end{enumerate}

\section{Dataset MedQuAD}

O \textbf{MedQuAD} (\textit{Medical Question Answering Dataset}) é o dataset base utilizado para o fine-tuning do modelo Mistral 7B. Criado pelo National Institutes of Health (NIH) dos Estados Unidos, o MedQuAD reúne pares de perguntas e respostas médicas extraídas de fontes confiáveis do governo americano.

\subsection{Escolha da Versão e Limitações de Copyright}

Existem múltiplas versões do MedQuAD disponíveis no HuggingFace. Optou-se pela versão \texttt{lavita/MedQuAD}, que possui os metadados mais completos. No entanto, uma \textbf{limitação importante} foi descoberta durante a implementação:

\begin{warningbox}[Remoção de Respostas por Copyright]
O dataset \texttt{lavita/MedQuAD} contém 47.441 registros, porém \textbf{aproximadamente 65\% (31.034 registros) possuem o campo \texttt{answer=None}}. Isso ocorreu porque as instituições originais (NIH, NCI, etc.) solicitaram a remoção das respostas por questões de copyright, mantendo apenas as perguntas. O dataset utilizável contém \textbf{aproximadamente 16.400 pares QA completos}.
\end{warningbox}

\begin{table}[H]
\centering
\caption{Comparação de versões do MedQuAD disponíveis no HuggingFace.}
\label{tab:versoes-medquad}
\rowcolors{2}{fiap-light}{white}
\begin{tabularx}{\textwidth}{l c c X}
\toprule
\textbf{Dataset} & \textbf{Registros} & \textbf{Com Resposta} & \textbf{Observações} \\
\midrule
\texttt{lavita/MedQuAD} & 47.441 & \textbf{$\sim$16.400} & Metadados ricos, 65\% sem resposta \\
\texttt{keivalya/MedQuad-...} & 16.400 & 16.400 & Versão já filtrada \\
\texttt{Tonic/medquad} & 15.549 & 15.549 & Pré-formatado Alpaca \\
\bottomrule
\end{tabularx}
\end{table}

\subsection{Características do Dataset}

\begin{table}[H]
\centering
\caption{Resumo do dataset MedQuAD (versão lavita, registros válidos).}
\label{tab:medquad}
\rowcolors{2}{fiap-light}{white}
\begin{tabularx}{\textwidth}{l X}
\toprule
\textbf{Característica} & \textbf{Descrição} \\
\midrule
Tamanho original  & 47.441 registros (65\% com \texttt{answer=None}) \\
Tamanho útil      & \textbf{$\sim$16.400} pares pergunta-resposta completos \\
Idioma original   & Inglês \\
Fontes            & 8 institutos NIH: NCI, GHR, MedlinePlus, NIDDK, NINDS, NIHSeniorHealth, NHLBI, CDC \\
Metadados         & Códigos UMLS, tipos semânticos, sinônimos, URLs das fontes \\
Domínios          & Doenças, tratamentos, medicamentos, prevenção, diagnósticos \\
Qualidade         & Alta --- respostas baseadas em fontes governamentais revisadas \\
Licença           & Creative Commons BY 4.0 \\
\bottomrule
\end{tabularx}
\end{table}

\subsection{Abordagem Híbrida de Treinamento}

Com base em pesquisa sobre fine-tuning multilíngue, adotamos uma \textbf{abordagem híbrida} para composição do dataset final, combinando três fontes:

\begin{enumerate}
    \item \textbf{MedQuAD traduzido (80\%):} Pares QA traduzidos para português via Google Translate;
    \item \textbf{MedQuAD original em inglês (20\%):} Amostra dos dados originais para preservar a qualidade médica e a capacidade multilíngue do modelo;
    \item \textbf{QA de medicina brasileira:} Perguntas e respostas nativas em português sobre temas específicos do Brasil.
\end{enumerate}

\begin{infobox}[Fundamentação da Abordagem Híbrida]
Pesquisas recentes (2024-2025) demonstram que fine-tuning exclusivamente com dados traduzidos pode introduzir ruído e perder nuances do conhecimento médico original. A inclusão de 20\% de dados em inglês, combinada com dados nativos em português, mitiga esses problemas e preserva a capacidade multilíngue do Mistral 7B.
\end{infobox}

\subsection{QA de Medicina Brasileira}

Para cobrir conhecimento médico específico do Brasil --- ausente no MedQuAD americano ---, desenvolvemos um conjunto de \textbf{55 pares QA} sobre temas brasileiros:

\begin{table}[H]
\centering
\caption{Categorias de QA de medicina brasileira.}
\label{tab:qa-brasileiro}
\rowcolors{2}{fiap-light}{white}
\begin{tabularx}{\textwidth}{l c X}
\toprule
\textbf{Categoria} & \textbf{QA} & \textbf{Exemplos de Tópicos} \\
\midrule
Arboviroses & 6 & Dengue, Zika, Chikungunya, diagnóstico diferencial \\
Malária & 5 & Regiões endêmicas, P. vivax vs P. falciparum, tratamento \\
Doença de Chagas & 4 & Transmissão oral (açaí), benznidazol, cardiopatia \\
Leishmaniose & 3 & Tegumentar vs visceral, antimoniato de meglumina \\
Febre Amarela & 2 & Vacinação, áreas de recomendação \\
Acidentes com animais & 7 & Jararaca, cascavel, aranha-marrom, escorpião \\
SUS e navegação & 5 & UBS, CAPS, referência/contrarreferência, alto custo \\
Medicamentos Brasil & 4 & Dipirona, RENAME, Farmácia Popular \\
Vacinação PNI & 3 & Calendário infantil, febre amarela, dengue \\
Pré-natal SUS & 3 & Consultas, exames obrigatórios, diabetes gestacional \\
Saúde Mental & 2 & RAPS, tipos de CAPS, manejo de crise \\
\midrule
\textbf{Total} & \textbf{55} & \\
\bottomrule
\end{tabularx}
\end{table}

Estes QA foram elaborados com base em fontes oficiais: Ministério da Saúde, Protocolos Clínicos e Diretrizes Terapêuticas (PCDT), RENAME 2024 e diretrizes do PNI.

\subsection{Pipeline de Preparação}

O \texttt{scripts/03\_preparar\_dataset.py} executa o seguinte pipeline:

\begin{enumerate}
    \item \textbf{Download:} Baixa o dataset \texttt{lavita/MedQuAD} do HuggingFace;
    \item \textbf{Filtragem:} Remove registros com \texttt{answer=None} ($\sim$65\% do total);
    \item \textbf{Tradução:} Traduz perguntas e respostas para português (com checkpoints a cada 1.000 registros);
    \item \textbf{Amostragem EN:} Seleciona 20\% do MedQuAD original em inglês;
    \item \textbf{QA Brasileiro:} Carrega os 55 pares de medicina brasileira;
    \item \textbf{Composição:} Combina as três fontes e embaralha;
    \item \textbf{Formatação Alpaca:} Salva no formato JSONL.
\end{enumerate}

O formato Alpaca utilizado segue a estrutura padrão para instruction tuning:

\begin{lstlisting}[caption={Formato Alpaca para fine-tuning.}]
{
    "instruction": "Responda a seguinte pergunta medica de forma
                    clara e detalhada.",
    "input": "Quais sao os sintomas do diabetes tipo 2?",
    "output": "Os sintomas comuns do diabetes tipo 2 incluem..."
}
\end{lstlisting}

\subsection{Composição Final do Dataset}

\begin{table}[H]
\centering
\caption{Composição do dataset final de treinamento.}
\label{tab:composicao-dataset}
\rowcolors{2}{fiap-light}{white}
\begin{tabular}{l r r l}
\toprule
\textbf{Fonte} & \textbf{Registros} & \textbf{Proporção} & \textbf{Arquivo} \\
\midrule
MedQuAD traduzido (PT) & $\sim$16.400 & 78\% & \texttt{medquad\_pt.jsonl} \\
MedQuAD original (EN) & $\sim$4.100 & 20\% & \texttt{medquad\_en\_sample.jsonl} \\
QA Brasileiro & 55 & 2\% & \texttt{brazilian\_medical\_qa.jsonl} \\
\midrule
\textbf{Total} & \textbf{$\sim$20.500} & 100\% & \texttt{training\_data.jsonl} \\
\bottomrule
\end{tabular}
\end{table}

\begin{infobox}[Arquivos Versionados]
Os arquivos de dataset processado (\texttt{medquad\_pt.jsonl}, \texttt{medquad\_en\_sample.jsonl}, \texttt{brazilian\_medical\_qa.jsonl} e \texttt{training\_data.jsonl}) são versionados no repositório Git. Isso permite reproduzir o treinamento sem necessidade de re-executar a tradução, que leva várias horas.
\end{infobox}
